% Cover letter - Ryan Full Stack AI Engineer
\documentclass[letterpaper,11pt]{article}
\usepackage[empty]{fullpage}
\usepackage[hidelinks]{hyperref}
\usepackage[english]{babel}
\usepackage{parskip}
\setlength{\parskip}{6pt}
\usepackage{setspace}
\addtolength{\oddsidemargin}{-0.5in}
\addtolength{\evensidemargin}{-0.5in}
\addtolength{\textwidth}{0.95in}
\addtolength{\topmargin}{-.80in}
\addtolength{\textheight}{1.7in}
\urlstyle{same}
\raggedright
\setstretch{1.0}
\pagestyle{empty}
\begin{document}

Dear Hiring Manager,

I am writing to apply for the Full Stack AI Engineer position at Ryan. I am currently completing an MSc in Health Data Science at University College London, where a distinction is expected, after graduating from Imperial College London. Across my academic and professional experience, I have built applications that use large language models, retrieval and agentic workf\kern0pt lows to solve problems that people actually have at work. I am particularly interested in this role because it combines software engineering with the judgement of deciding where generative AI genuinely improves a process and where ordinary deterministic software is the better answer.

What attracts me most to Ryan is the opportunity to work on AI-enabled applications end to end, from understanding a business workf\kern0pt low through to development, deployment and iteration with real users. I enjoy taking an unclear business problem, identifying what users actually need, and turning it into a workf\kern0pt low that can be tested and improved. The posting's emphasis on engineering fundamentals and learning agility over knowledge of any particular framework describes how I have worked so far.

At TikTok Shop UK I am a Data Analyst Intern in Business Planning and Operations. Alongside my analytical work I built an AI agent that automates repetitive steps of the analytics workf\kern0pt low, which has been adopted by my own team and by partner teams. Building it meant understanding the existing workf\kern0pt low in detail, deciding which steps were worth automating and which needed a person, and then supporting other people as they started to use it. I use Claude Code as part of my development work and review the code it produces rather than accepting it unchecked.

Before TikTok I worked as a Data Scientist Intern at DS Digital Technology on a medical AI engagement for a pharmaceutical client. I developed a Python application that called six LLM APIs with retrieval and search augmentation, tool use and multi-step workf\kern0pt lows, generating and scoring over 2,400 responses. I designed a 246-question benchmark and a six-dimension rubric with red-line rules, which the client adopted as its acceptance standard, and co-built a medically reviewed knowledge graph of 123 entities and 161 relations used for retrieval. I then ran blinded paired A/B tests to measure whether the knowledge graph actually improved answer quality, because I did not want to assume that adding retrieval had helped.

My MSc dissertation is a full pipeline rather than a single model: cardiovascular risk prediction on 458,840 UK Biobank participants and doubly robust causal models on a further cohort, delivered as a deployed Streamlit application so that people can run scenarios without touching the code. I also led a group project building a clinical NLP pipeline that parsed radiology reports into structured features using f\kern0pt ine-tuned transformer models. Both projects involved writing code that someone other than me had to run.

At Ipsos Strategy3 I worked directly with client stakeholders, turning open-ended business questions into structured analyses and presenting the results to people who wanted the recommendation rather than the method. This experience has strengthened my ability to explain technical decisions and trade-of\kern0pt fs to both technical and non-technical audiences, which I understand is part of this role.

I should be clear about where I am starting from. My full-stack work has been in Python, with Streamlit for interfaces, rather than in React and Node or TypeScript, and my cloud experience with A\kern0pt WS is at a basic level. I am comfortable with APIs, databases, Git, testing and debugging, I learn new frameworks quickly, and I would expect to be productive on a React and FastAPI stack early on with the mentorship the role describes. I would bring Ryan solid engineering fundamentals, genuine hands-on experience building and evaluating LLM and agentic applications, and the habit of checking whether what I built actually worked. I have the right to work in the UK until January 2029 without sponsorship.

Thank you for considering my application. I would welcome the opportunity to discuss how my experience building AI-enabled applications could contribute to your engineering team.

\vspace{2pt}
Yours sincerely,

\vspace{8pt}
Zhuoyan Sun

\end{document}
